\begin{figure}[H]
\centering
\resizebox{0.94\textwidth}{!}{%
\begin{tikzpicture}[
 panel/.style={draw=prgray!55,rounded corners=3pt,minimum width=52mm,minimum height=42mm},
 head/.style={font=\small\bfseries,align=center},
 form/.style={draw,rounded corners=2pt,fill=prlightblue,minimum width=19mm,minimum height=8mm,align=center,font=\scriptsize},
 mini/.style={draw,rounded corners=2pt,fill=prlightblue,minimum height=7mm,align=center,font=\scriptsize,inner xsep=1.5mm},
 op/.style={circle,draw,fill=prlightgold,minimum size=6mm,font=\scriptsize,inner sep=0pt},
 place/.style={circle,draw,minimum size=4.2mm,inner sep=0pt},
 trans/.style={rectangle,draw,fill=prlightteal,minimum width=7mm,minimum height=4.2mm,font=\scriptsize},
 tau/.style={rectangle,draw,fill=prgray!20,minimum width=1.5mm,minimum height=6mm,inner sep=0pt},
 arr/.style={-{Stealth[length=1.4mm]},thin}
]
\node[panel] (p1) at (0,0) {};
\node[head] at (0,1.7) {Duplicate prefix / silent routing};
\node[form] (dleft) at (-1.45,0.45) {$A_0A_1\ \mid\ A_0A_2$};
\node[font=\small] at (0,0.45) {$\Longleftrightarrow$};
\node[form] (dright) at (1.45,0.45) {$A_0\,\tau\,\{A_1,A_2\}$};
\node[font=\scriptsize,align=center] at (0,-1.75) {same visible traces;\\different internal routing};

\node[panel] (p2) at (5.6,0) {};
\node[head] at (5.6,1.7) {Concurrency / interleaving};
\node[form] at (4.2,0.45) {$\mathsf{AND}(A_0,A_1)$};
\node[font=\small] at (5.6,0.45) {$\Longleftrightarrow$};
\node[form] at (7.0,0.45) {$A_0A_1\ \mid\ A_1A_0$};
\node[font=\scriptsize,align=center] at (5.6,-1.75) {same complete traces;\\different partial order};

\node[panel] (p3) at (11.2,0) {};
\node[head] at (11.2,1.7) {Block structure / M-pattern};
\node[mini] at (11.2,1.05) {$\mathsf{XOR}(A_0,\mathsf{AND}(A_1,A_2))$};
\node[font=\small] at (11.2,0.55) {$\Updownarrow$};
% Minimal M-pattern fragment: A0 shares both input places, whereas A1 and A2
% each share only one.  Omitting the surrounding source/sink keeps the
% non-free-choice relation legible at column width.
\node[place] (mp1) at (10.30,-0.18) {};
\node[place] (mp2) at (10.30,-0.82) {};
\node[trans] (m1) at (11.20,0.15) {$A_1$};
\node[trans,fill=prlightred] (mA) at (11.20,-0.50) {$A_0$};
\node[trans] (m2) at (11.20,-1.15) {$A_2$};
\node[place] (mq1) at (12.10,0.15) {};
\node[place] (mqA) at (12.10,-0.50) {};
\node[place] (mq2) at (12.10,-1.15) {};
\draw[arr] (mp1)--(m1); \draw[arr] (m1)--(mq1);
\draw[arr] (mp1)--(mA); \draw[arr] (mp2)--(mA); \draw[arr] (mA)--(mqA);
\draw[arr] (mp2)--(m2); \draw[arr] (m2)--(mq2);
\node[font=\scriptsize,align=center] at (11.2,-1.75) {same complete traces;\\non-free-choice routing};
\end{tikzpicture}%
}
\caption{Controlled representation contrasts. Every panel uses the same activity notation. The double arrow denotes agreement in visible complete traces. Concurrency/interleaving deliberately changes partial-order evidence; the duplicate-prefix and M-pattern pairs retain the stronger recorded relation.}
\label{fig:motifs}
\end{figure}
